% Rúbrica: Análisis exploratorio y modelación (la variable objetivo y su métrica).
\section{Definición del objetivo}\label{sec:objetivo}

El modelo aprende a ordenar juegos según una señal que sale de sus reseñas: un voto negativo
escrito antes de las dos horas de juego. Esa señal no es el arrepentimiento mismo, el corte de
120 minutos viene de una regla de Steam y la métrica tiene que funcionar con una clase muy rara.

\subsection{La fórmula}\label{sec:objetivo-formula}

Para cada reseña $i$,
\[
  Y_i =
  \begin{cases}
    1 & \text{si } \texttt{playtime\_at\_review}_i < 120 \text{ y } \texttt{voted\_up}_i = 0,\\
    0 & \text{en otro caso,}
  \end{cases}
\]
donde \texttt{playtime\_at\_review} son los minutos que llevaba jugados quien escribió la reseña
cuando la publicó y \texttt{voted\_up} es el pulgar: 1 arriba, 0 abajo. En data-v1 hay
\cifraPositivosEntrenamiento{} reseñas con $Y = 1$ de \cifraResenasEntrenamiento{}, una
prevalencia de \cifraPrevalenciaEntrenamientoPorResena{} por reseña. A esa etiqueta se le llama
\emph{señal de arrepentimiento temprano}.

\subsection{Por qué es una señal proxy}\label{sec:objetivo-proxy}

Steam no registra arrepentimiento. Registra que alguien votó en contra y cuánto había jugado al
escribir. Los 120 minutos vienen de la política de reembolsos de la tienda: se puede pedir el
reembolso dentro de las dos semanas posteriores a la compra y con menos de dos horas de juego
\parencite{steam_reembolsos}. Una reseña negativa escrita dentro de esa ventana es lo más cercano
que dejan los datos a «esto no era para mí, y me di cuenta pronto».

Por eso el sitio y este documento la llaman señal de arrepentimiento temprano y nunca dicen que
alguien se arrepintió. La señal deja fuera a quien se arrepiente y no escribe, a quien pide el
reembolso sin reseñar y a quien reseña después de las dos horas. Tampoco mide abandono ni
reembolsos: de ninguno de los dos hay dato.

Lo que sí muestran los datos es que las negativas llegan antes. El
\cifraNegativasAntesDelUmbralPorResena{} de las reseñas negativas se escribe antes de los 120
minutos, contra el \cifraPositivasAntesDelUmbralPorResena{} de las positivas. La señal se queda con
ese tramo temprano de las negativas.

\subsection{Robustez del umbral de 120 minutos}\label{sec:objetivo-umbral}

El umbral no se eligió con los datos, sale de la regla de Steam. Lo que sí se revisó es si
moverlo cambia qué juegos tienen más señal. Con 60, 90 o 180 minutos cambia cuántas reseñas
cuentan, pero el orden de los juegos casi no se mueve: su correlación de Spearman con el orden de
120 va de \cifraSpearmanUmbralMin{} a \cifraSpearmanUmbralMax{} (tabla~\ref{tab:umbral}).

\begin{table}[htbp]
  \centering\small
  % Generado por documento/generar_figuras.py. No se edita a mano.
\begin{tabular}{rrrr}
\toprule
Umbral (min) & Reseñas con señal & Prevalencia por reseña & Spearman con 120 (por juego) \\
\midrule
60 & 1,649 & 1.33\,\% & 0.964 \\
90 & 2,250 & 1.81\,\% & 0.989 \\
120 & 2,714 & 2.19\,\% & 1.000 \\
180 & 3,349 & 2.70\,\% & 0.993 \\
\bottomrule
\end{tabular}

  \caption{La señal con otros umbrales. La prevalencia es por reseña; el Spearman compara las
  tasas por juego de los \cifraJuegosEntrenamiento{} juegos contra las de 120 minutos. Fuente:
  data-v1, con \texttt{sensibilidad\_umbral} de \texttt{analisis/exploracion.py}.}
  \label{tab:umbral}
\end{table}

La figura~\ref{fig:minutos} muestra en qué minuto se escribe cada reseña. En 120 no hay escalón:
por minuto, las negativas pasan de \cifraNegativasPorMinutoAntesDelUmbral{} a
\cifraNegativasPorMinutoDespuesDelUmbral{} al cruzarlo, y las positivas de
\cifraPositivasPorMinutoAntesDelUmbral{} a \cifraPositivasPorMinutoDespuesDelUmbral{}. El salto
que se ve está en los 180 minutos y solo en las positivas, que pasan de
\cifraPositivasPorMinutoAntesDeTresHoras{} a \cifraPositivasPorMinutoDesdeTresHoras{} reseñas por
minuto en \cifraJuegosSaltoTresHoras{} juegos. Las negativas apenas cambian (de
\cifraNegativasPorMinutoAntesDeTresHoras{} a \cifraNegativasPorMinutoDesdeTresHoras{}). La causa de
ese salto no está en estos datos, y no toca la señal, que se define en 120 minutos y con reseñas
negativas.

\begin{figure}[htbp]
  \centering
  \includegraphics[width=\textwidth]{minutos-al-resenar.pdf}
  \caption{Minutos jugados al escribir la reseña, por pulgar, en escala logarítmica. Cada curva
  se normaliza sobre su propio grupo. Los intervalos tienen el mismo ancho en escala log, con un borde en 180
  y ninguno en 120. Fuente: data-v1.}
  \label{fig:minutos}
\end{figure}

\subsection{Por qué PR-AUC}\label{sec:objetivo-metrica}

Con una prevalencia de \cifraPrevalenciaEntrenamientoPorResena{}, un modelo que siempre dice «no»
acierta el \cifraExactitudTrivialPorResena{} de las veces y no encuentra un solo caso. La
exactitud premia justo eso, así que aquí no sirve. El proyecto usa el área bajo la curva de
precisión y exhaustividad (PR-AUC, con \texttt{average\_precision\_score} de scikit-learn). Mide
qué proporción de lo que el modelo pone arriba es de verdad señal, y con clases desbalanceadas informa
más que la curva ROC \parencite{saito2015,davis2006}.

El PR-AUC de un clasificador que no sabe nada es la prevalencia, \cifraPRAUCTrivial{}. Ese es el
piso, y por eso en este documento el PR-AUC del modelo siempre va junto al del trivial. En la
validación agrupada por juego el modelo llega a \cifraPRAUCModelo{}, \cifraPRAUCCociente{} veces
ese piso (sección~\ref{sec:evaluacion}).
